\subsection*{5. A hue and cry against so-called ``general'' topology, and heuristic reflections toward a so-called ``tame'' topology}
The fifth section starts on p.~25 with the problem of reconstructing a
stratified structure from canonical ``pièces de construction'', and
retraces a change from the successive singular loci of a variety to the
stratifications arising from degeneration in moduli problems. The
Mumford-Deligne compactifications on p.~26 connect this problem with the
Teichmüller tower of \S2. The two-stratum model uses a tubular neighborhood
$T$ of $Y$ in $X$, its complement $U=X\setminus\mathrm{Int}(T)$, and a
gluing map $f:\partial U\longrightarrow Y$; p.~27 reconstructs the space
by the stated amalgamated sum. The description on p.~26 contains a
mathematical slip in the base of the associated cone bundle: the preceding
bundle has projection $p:\partial T\longrightarrow Y$, and its fiberwise
cone is explicitly
$((\partial T\times[0,1])\amalg Y)/((z,0)\sim p(z))$,
with projection $[z,t]\mapsto p(z)$ and $y\mapsto y$ to $Y$.
This formula exhibits the base of that construction; the printed phrase
``fibré en cônes sur $T$'' remains unchanged in the readers. On pp.~27--28
the lack of canonical choices leads to ``catégories isotopiques'',
resuming Ladegaillerie's work from \S3, and then to the author's account
of the limits of his initial hopes in the presence of wild phenomena.
Hironaka's work prompts the renewed search for a ``topologie modérée''
(pp.~28--29). The polemic on pp.~29--30 explains the geometric
requirements behind that search. From p.~30 onward Grothendieck prefers
an axiomatic treatment encompassing several tame theories, with
comparison maps between spaces of isomorphisms and their homotopy types
(p.~31); triangulability is the strongest axiom in his provisional system,
and he explicitly records the unproved state of the proposed comparison
theorem (p.~32). The discussion of fields and point dimensions on
pp.~32--33 explains one advantage he sees in retaining several theories.
Another mathematical slip occurs in the collar description on p.~33,
where the product is written with $T$ as its factor. The elementary collar
of $T=[-1,1]$ is the map
$c:\{-1,1\}\times[0,\varepsilon)\longrightarrow[-1,1]$,
$c(s,t)=s(1-t)$, with $0<\varepsilon<1$; its inverse on the two endpoint
neighborhoods is $x\mapsto(\operatorname{sgn}(x),1-|x|)$.
The displayed isomorphism uses $\partial T$ as the product factor;
the source's plain $T$ is preserved. On p.~34 he allows that the proposed
equivalence of isotopy categories may be too optimistic, while defending
the content of the dévissage already developed. Local retractions provide
additional structure and lead to spaces indexed by flags in the ordered
set of strata, but he immediately notes that these extra data are absent
from the geometric situations that motivated the work. On p.~35 the
moduli example replaces the index set by the category of ``graphes MD'',
and the argument returns to the combinatorial tower of \S2. The
``topossique'' test on pp.~35--36 carries the same heuristic into étale
topoi and the Galois-compatible dévissage of profinite Teichmüller groups;
Grothendieck proposes a ``voisinage tubulaire canonique d'un sous-topos''
and a comparison with the transcendental setting. The final change of
emphasis, on p.~36, is that direct geometric arguments can probably
short-circuit much of this theory for the moduli applications, although
he retains its value as a guide. The later note (6), pp.~55--56,
qualifies the polemic by recognizing the place and beauty of wild
phenomena; note (7), pp.~56--57, gives the morphism data and the hexagonal
diagram behind the flag-indexed construction. Related titles are
\href{https://grothendieck.umontpellier.fr/archives-grothendieck/}{\emph{Topologie modérée} (cotes 120--121), \emph{Espaces stratifiés et voisinages côniques} (cote 150), and \emph{Topos stratifié} (cote 151)},
and Lou van den Dries's
\href{https://assets.cambridge.org/97805215/98385/frontmatter/9780521598385_frontmatter.pdf}{\emph{Tame Topology and O-minimal Structures}}.
(\href{https://grothendieck.umontpellier.fr/web/viewer.html?file=../119.pdf\#page=28}{Copy 1, pp.~25--36, beginning at site 27};
\href{https://grothendieck.umontpellier.fr/web/viewer.html?file=../119.pdf\#page=29}{cone-bundle passage on p.~26, site 28};
\href{https://grothendieck.umontpellier.fr/web/viewer.html?file=../119.pdf\#page=36}{collar passage on p.~33, site 35};
\href{https://grothendieck.umontpellier.fr/web/viewer.html?file=../119.pdf\#page=39}{closing reassessment on p.~36, site 38}.)
